\bookchapter{Capitolo 6: Proteggi la tua vita digitale — Sicurezza e privacy}

\chapterintro{Sicurezza del telefono = privacy + patrimonio. In un'epoca in cui i dati sono un asset, proteggere il telefono significa proteggere te stesso.}

\booksection{Blocco schermo: la prima linea di difesa fisica}

Il blocco schermo è la prima linea di difesa del telefono e anche la più trascurata.

\booksubsection{Classifica per livello di sicurezza:}

\begin{enumerate}
  \item \textbf{Password/PIN}: la massima sicurezza, è la garanzia finale (quando dimentichi impronta/volto).
  \item \textbf{Impronta digitale}: il perfetto equilibrio tra sicurezza e comodità.
  \item \textbf{Riconoscimento facciale}: il più comodo, ma la sicurezza dipende dalla tecnologia (Face ID 3D sicuro, 2D aggirabile con foto).
  \item \textbf{Sequenza}: comoda ma lascia tracce sullo schermo, facile da indovinare.
\end{enumerate}

\booksubsection{Imposta un blocco schermo forte:}

\begin{itemize}
  \item Imposta obbligatoriamente una password complessa (non compleanno, non 123456).
  \item Attiva il riconoscimento biometrico.
  \item Imposta il blocco automatico su "immediato" o sul tempo più breve.
\end{itemize}

\booksection{“Trova il mio telefono”: lo strumento antifurto definitivo}

\textbf{Deve essere attivato!}Permette di:

\begin{itemize}
  \item \textbf{Localizzare}: dov'è il telefono, su una mappa.
  \item \textbf{Riprodurre un suono}: fa squillare il telefono, utile in casa.
  \item \textbf{Modalità smarrito}: blocca il telefono da remoto e mostra un numero di contatto.
  \item \textbf{Cancella da remoto}: l'ultima risorsa, elimina tutti i dati (ma dopo non è più possibile localizzarlo).
\end{itemize}

Vai in [Impostazioni]; su Apple cerca “Trova il mio iPhone”, su Android “Trova il mio dispositivo” o un servizio equivalente del brand; assicurati che sia attivo (verde).

\booksection{Aggiornamenti di sistema: applicare patch, tappare le falle}

Lo scopo principale degli aggiornamenti è correggere vulnerabilità di sicurezza, come rattoppare i buchi nel muro di casa. Quando vedi la notifica di aggiornamento, installa il prima possibile.

Naturalmente, gli aggiornamenti major (ad esempio da iOS 17 a iOS 18) possono avere problemi di compatibilità: puoi aspettare qualche giorno prima di installare.

\booksection{Riconoscere e prevenire le truffe telecomunicazioni}

Con lo sviluppo rapidissimo dell'AI, il modello cognitivo tradizionale “vedere per credere” è ormai superato. I truffatori sfruttano tecnologie AI come il deepfake per passare dalla truffa “a pioggia” al crimine “mirato”, con un salto di qualità totale.

\booksubsection{Le truffe più comuni nell'era dell'AI}

\textbf{(1) Truffa con face-swap e clonazione vocale AI}\par

I criminali raccolgono illegalmente foto ad alta risoluzione pubblicate sui social o campioni vocali di pochi secondi e, con il deepfake, clonano “video di amici/parenti” o “voci” quasi indistinguibili dal vero. Copioni frequenti:\textbf{incidente stradale, prestito urgente, rapimento}. Queste truffe sfruttano legami affettivi e\textbf{sono molto ingannevoli}.

\textbf{(2) Falsificazione di autorità / servizio clienti}\par

L'AI può simulare con precisione voce e tono di addetti al servizio clienti ufficiali o di funzionari di polizia/giustizia, e persino comporre screenshot di bonifici, documenti ufficiali, mandati d'arresto altamente realistici;\textbf{creando panico o pressione autoritaria}, inducono la vittima a trasferire denaro o ad attivare la condivisione schermo, sottraendo informazioni sull'account.

\textbf{(3) Truffa di investimenti AI}\par

Tramite algoritmi AI vengono falsificati “screenshot di guadagni sicuri”, interfacce di software di trading quantistico o “report di analisi da maestro”, adescando con alti rendimenti, vendendo servizi di investimento costosi o spingendo la vittima a trasferire fondi su conti di trading controllati dai criminali, per poi\textbf{scaricare tutto e sparire}.

\booksubsection{Come difendersi — la logica dello “specchio rivelatore”}

L'AI può clonare voce e volto, ma non è invincibile: ci sono alcune crepe da sfruttare.

\begin{itemize}
  \item \textbf{Verifica con azioni}: in videochiamata, chiedi esplicitamente all'interlocutore di compiere azioni coordinate predefinite, come "passa la mano velocemente davanti al viso", "gira il volto per mostrare il contorno dell'orecchio" o "ammicca rapidamente 5 volte". Negli scenari con occlusioni dinamiche complesse, movimenti rapidi e variazioni di luce, il face-swap AI\textbf{mostra facilmente difetti come scattosità, distorsione del volto o sbavatura dei bordi}.
  \item \textbf{Verifica con codici privati}: stabilisci e usa domande di verifica private predefinite, come "cosa hai ordinato l'ultima volta che siamo andati a cena insieme", "qual è stato il primo regalo che ti ho fatto", dettagli noti solo a voi due. Se l'altra parte esita o risponde a sproposito,\textbf{interrompi subito la chiamata e attiva il piano di sicurezza}.
  \item \textbf{Doppia verifica}: segui rigorosamente il principio del Ministero della Pubblica Sicurezza "\textbf{quattro no e una verifica}" — non credere a messaggi sconosciuti, non trasferire denaro a cuor leggero, non divulgare dati personali sensibili, non diffondere contenuti sospetti. In caso di movimento di denaro,\textbf{richiama sempre il numero originale già noto o verifica di persona}. Per i bonifici, scegli la modalità "accredito in 24 ore", per lasciare tempo sufficiente a individuare la truffa e recuperare i fondi.
\end{itemize}

\booksubsection{Truffe tradizionali (da tenere comunque d'occhio):}

\begin{itemize}
  \item \textbf{Vincita / cashback per "missioni"}: adescamento, prima ti danno un piccolo guadagno poi ti truffano cifre grosse.
  \item \textbf{Truffa sentimentale online}: "pig butchering", coltiva il sentimento e poi truffa il denaro.
  \item \textbf{Falsi conoscenti}: rubano l'account e chiedono prestiti (più difficile da distinguere nell'era AI).
\end{itemize}

\booksubsection{Protezioni integrate nel telefono:}

\begin{itemize}
  \item Blocco delle chiamate moleste.
  \item Protezione dal furto di codici di verifica.
  \item Controllo di sicurezza dei link negli SMS.
\end{itemize}

\booksubsection{motto antitruffa:}

\textbf{Non credere, non divulgare, non cliccare, non trasferire, in caso di dubbio verifica prima!}Davanti a un'“emergenza” prenditi 5 minuti per calmarti: quasi tutte le truffe si tradiscono in quel momento.

\booksection{Difendersi dal “doxxing” e dagli attacchi di social engineering: spezza la catena della perdita di privacy}

Il principio fondamentale del “doxxing” (attacco di social engineering) consiste nel\textbf{raccolta sistematica, analisi e assemblaggio di frammenti digitali personali per costruire un profilo completo dell'individuo}. La chiave per difendersi sta nel tracciare confini chiari per la privacy personale e stabilire un sistema di protezione multilivello.

\booksubsection{Anonimizzazione delle foto}

I metodi base per la pulizia EXIF sono descritti nel Capitolo 2 “Gestione di immagini e file system”. Prima di condividere foto sui social media, è necessario\textbf{eliminare definitivamente i metadati della foto (informazioni EXIF)}, inclusi data di scatto, modello del dispositivo, coordinate GPS e altri dati sensibili. Le informazioni chiave presenti nelle foto, come\textbf{numeri civici, targhe, codici a barre dei biglietti, numeri di documento d'identità}, devono essere rigorosamente offuscate o sfocate per evitare di rivelare inconsapevolmente le tracce dei propri spostamenti.

\begin{itemize}
  \item iOS: in [Impostazioni] > [Privacy e sicurezza] > [Fotocamera] è possibile disattivare "Posizione precisa"; durante la condivisione selezionare "Solo foto corrente" per rimuovere i metadati.
  \item Android: alcune app di galleria (come Google Foto) chiedono se mantenere le informazioni sulla posizione durante la condivisione.
  \item Strumenti universali: ExifTool, Metapho ecc. per visualizzare/modificare EXIF.
\end{itemize}

\booksubsection{Restringimento dei permessi delle app}

I metodi base di gestione dei permessi sono descritti nel dettaglio nel Capitolo 4 “Gestione dei permessi”. Nello scenario “anti-doxxing”, è inoltre necessario revisionare periodicamente e disattivare i permessi non essenziali di app di e-commerce e strumenti, come\textbf{microfono, rubrica, posizione}, per evitare che le app “ascoltino” di nascosto le conversazioni quotidiane in background o raccolgano illegalmente reti di relazioni personali; le app a lungo inutilizzate ma ancora dotate di permessi elevati dovrebbero essere disinstallate tempestivamente, riducendo alla radice i punti di rischio per la perdita di privacy.

\begin{itemize}
  \item Quando vedi una richiesta di "permesso rubrica", chiediti prima: quest'app ha davvero bisogno di conoscere tutti i miei contatti?
  \item App non usate da tre mesi, disinstallale senza esitazione. I permessi si restringono con l'uso, non si allargano.
\end{itemize}

\booksubsection{Messa in sicurezza dei dispositivi smart}

Telecamere smart, serrature smart, casse smart e altri dispositivi IoT al primo utilizzo devono\textbf{avere immediatamente la password di accesso predefinita modificata}, impostando una password indipendente complessa con lettere maiuscole e minuscole, numeri e caratteri speciali; le funzioni di accesso remoto non utilizzate di frequente dovrebbero restare disattivate, controllando regolarmente gli aggiornamenti firmware dei dispositivi e correggendo tempestivamente le vulnerabilità di sicurezza,\textbf{per evitare che hacker si introducano e spiino lo spazio privato domestico}.

\begin{tcbwarn}{Avvertenza da caso reale}
Nel 2022, la polizia di una località ha riferito che hacker, decifrando la password predefinita di telecamere domestiche di una certa marca (molti utenti non l'avevano mai modificata), hanno sottratto immagini delle camere da letto di migliaia di famiglie e le hanno vendute online.\textbf{Non cambiare la password predefinita = lasciare la porta di casa spalancata}.

\end{tcbwarn}

\booksection{Protezione della privacy: i tuoi dati sono la tua proprietà}

Le foto, le chat, le tracce di posizione nel telefono sono essenzialmente i tuoi “beni digitali”. Una volta trapelati, nel migliore dei casi fastidi pubblicitari, nel peggiore dei casi doxxing, truffe o persino ricatti.

\booksubsection{“Principio di minimizzazione” dei permessi delle app:}

\begin{itemize}
  \item \textbf{Posizione}: per taxi/mappe usa "Consenti sempre"; per le altre app usa "Consenti solo durante l'uso".
  \item \textbf{Rubrica}: concedila solo se necessario (come aggiungere amici su WeChat), altrimenti vietala sempre.
  \item \textbf{Galleria / Tutte le foto}: da iOS 14+ si può autorizzare "Foto selezionate", anche Android 13+ supporta l'autorizzazione granulare.
  \item \textbf{Microfono / Fotocamera}: tienili disattivati normalmente, attivali quando servono. Se sospetti di essere registrato di nascosto, scorri il Centro di controllo per vedere se c'è un "puntino verde/arancio".
  \item \textbf{Tracciamento in background}: da iOS 14.5+ è obbligatorio chiedere prima che le app usino IDFA; su Android si può disattivare "ID pubblicità".
\end{itemize}

\booksubsection{Pratica della gestione della privacy:}

\begin{itemize}
  \item \textbf{Rapporto privacy (iOS)}: [Impostazioni] > [Privacy e sicurezza] > [Rapporto privacy app], per vedere quali app accedono di nascosto a posizione, fotocamera, microfono.
  \item \textbf{Cronologia utilizzo permessi (Android)}: [Impostazioni] > [Privacy] > [Cronologia utilizzo permessi], per vedere chi ha richiamato permessi sensibili e quando.
  \item \textbf{Limita tracciamento pubblicitario}: disattiva "ID pubblicità" o "Limita tracciamento annunci" nelle impostazioni di sistema, per ridurre le "pubblicità mirate fastidiose".
\end{itemize}

\booksection{Lista di “primo soccorso” in caso di smarrimento / furto}

Se hai perso il telefono non farti prendere dal panico, segui quest'ordine:

\begin{enumerate}
  \item \textbf{Blocca immediatamente la SIM}: contatta l'operatore per sospenderla, evitando che venga usata per ricevere codici di verifica.
  \item \textbf{Blocco remoto}: tramite servizio cloud (iCloud / cloud del marchio) attiva "Modalità smarrimento".
  \item \textbf{Modifica le password importanti}: cambia subito le password di tutti gli account importanti (pagamenti, banca, email) che hai usato sul telefono.
  \item \textbf{Denuncia}: sporgi denuncia alla polizia, annotando il numero IMEI. Sarà utile anche per provare l'identità se il telefono viene ritrovato.
  \item \textbf{Avvisa parenti e amici}: per evitare che i truffatori chiedano prestiti a tuo nome.
  \item \textbf{(Per ultimo) Cancellazione remota}: esegui solo quando non ci sono più speranze di ritrovarlo; dopo la cancellazione la localizzazione non funziona più.
\end{enumerate}

\booksection{Sicurezza SIM ed eSIM}

\begin{itemize}
  \item \textbf{Imposta PIN della SIM}: una volta attivato, a ogni riavvio del telefono o cambio scheda sarà necessario inserire il PIN, per evitare che "si tolga la SIM e si cambi telefono" aggirando il blocco schermo.
  \item \textbf{Vantaggi eSIM}: non essendoci una scheda fisica estraibile, la sicurezza è maggiore; supporta inoltre la commutazione tra più numeri, molto comodo in viaggio.
  \item \textbf{Rischio di sostituzione secondaria della SIM}: attenzione alla truffa della "sostituzione secondaria operatore" - i truffatori usano la tua identità per rifare una SIM, prendendo così il controllo del tuo numero. Impostare "la sostituzione della SIM richiede la presenza personale del titolare con documento d'identità allo sportello" è la difesa più efficace.
\end{itemize}

\begin{tcbwarn}{Cose da non fare mai}
\begin{itemize}
  \item Usare password deboli come compleanno / 123456 / qwerty;
  \item Inviare il codice di verifica a chiunque (anche al "servizio clienti");
  \item Cliccare su link sospetti negli SMS, in particolare da numeri internazionali che iniziano con 00;
  \item Accedere alla banca online o ai pagamenti sotto Wi-Fi pubblico;
  \item Attivare "condivisione schermo" o "controllo remoto" con sconosciuti.
\end{itemize}

\end{tcbwarn}

\begin{tcbpractice}{Prova ad attivare e simulare il "blocco remoto"}
\begin{enumerate}
  \item Assicurati che la funzione "Trova il mio telefono" sia attiva.
  \item Su un browser di computer o su un altro telefono, accedi al corrispondente account di servizio cloud (iCloud.com o sito del servizio cloud del marchio).
  \item Trova la funzione "Trova dispositivo", prova a inviare al tuo telefono un comando "Riproduci suono" o simula un "Blocco", per familiarizzare con l'intera procedura.
\end{enumerate}

\end{tcbpractice}
